% !TEX program = pdflatex
% Анализ, верификация и сопоставление реализации непараметрических
% CUSUM-карт (по знакам SN и знаковым рангам SR) со средствами языка R.
% Первоисточники: Chakraborti & Graham (2019); Bakir & Reynolds (1979);
% Amin, Reynolds & Bakir (1995); Nascimento et al. (2026).
% Сборка: pdflatex (T2A + babel russian).
\documentclass[11pt,a4paper]{article}

\usepackage[T2A]{fontenc}
\usepackage[utf8]{inputenc}
\usepackage[russian]{babel}
\usepackage[margin=2.4cm]{geometry}
\usepackage{amsmath,amssymb}
\usepackage{pgfplots}
\pgfplotsset{compat=1.17}
\usepackage{booktabs}
\usepackage{array}
\usepackage{graphicx}
\usepackage{xcolor}
\usepackage{listings}
\usepackage{caption}
\usepackage{enumitem}
\usepackage[hidelinks]{hyperref}

\definecolor{accent}{RGB}{46,134,171}   % ACCENT  #2E86AB
\definecolor{alert}{RGB}{178,34,34}     % ALERT   #B22222
\definecolor{third}{RGB}{60,110,71}     % PALETTE #3C6E47
\definecolor{softgray}{RGB}{128,128,128} % CENTER #808080

\lstset{
  basicstyle=\ttfamily\small,
  commentstyle=\itshape\color{softgray},
  keywordstyle=\bfseries,
  breaklines=true,
  frame=single,
  rulecolor=\color{softgray},
  columns=fullflexible,
  keepspaces=true,
}

\captionsetup{font=small,labelfont=bf}
\setlist[itemize]{itemsep=1pt,topsep=2pt}
\setlist[description]{itemsep=1pt,topsep=2pt}
\emergencystretch=1.5em

\newcommand{\SN}{\mathrm{SN}}
\newcommand{\SR}{\mathrm{SR}}
\newcommand{\ARL}{\mathrm{ARL}}
\newcommand{\SDRL}{\mathrm{SDRL}}
\newcommand{\MRL}{\mathrm{MRL}}
\newcommand{\FAR}{\mathrm{FAR}}
\newcommand{\E}{\mathbb{E}}

\pgfplotsset{
  myaxis/.style={
    width=0.47\textwidth, height=0.36\textwidth,
    grid=both, grid style={softgray!25},
    axis line style={softgray}, tick align=outside,
    label style={font=\small}, tick label style={font=\small},
    legend style={font=\footnotesize, draw=softgray!50, fill=white},
    every axis plot/.append style={line width=1pt},
  },
}

\title{\bfseries Непараметрические CUSUM-карты по знакам (SN) и знаковым
рангам (SR):\\ математические основы, программная реализация,\\
верификация и сопоставление с языком R}
\author{}
\date{2026}

% Координаты графиков генерируются скриптом (не редактировать вручную):
%   poetry run python experiments/cusum/06_gen_figdata.py \
%       > docs/reports/cusum_macros.tex
% Сгенерировано experiments/cusum/06_gen_figdata.py.
% Не редактировать вручную: пересобрать командой make cusum.

% --- OC-кривые: CUSUM-SN, n=10, k=6, h=4 (рис. 4.5) ---
\def\ocSNNormal{(0.0000,464.984389) (0.2500,129.226695) (0.5000,25.643513) (0.7500,7.649991) (1.0000,3.513544) (1.2500,2.202467) (1.5000,1.654845) (1.7500,1.375193) (2.0000,1.213922)}
\def\ocSNT{(0.0000,464.984389) (0.2500,87.937635) (0.5000,14.801432) (0.7500,4.825277) (1.0000,2.615124) (1.2500,1.866863) (1.5000,1.527840) (1.7500,1.343294) (2.0000,1.231386)}
\def\ocSNLaplace{(0.0000,464.984389) (0.2500,52.675686) (0.5000,10.006446) (0.7500,4.072086) (1.0000,2.506005) (1.2500,1.888847) (1.5000,1.577720) (1.7500,1.394753) (2.0000,1.276433)}
\def\ocSNLogistic{(0.0000,464.984389) (0.2500,102.620136) (0.5000,18.373151) (0.7500,5.776767) (1.0000,2.948743) (1.2500,2.011903) (1.5000,1.597796) (1.7500,1.376067) (2.0000,1.242718)}
\def\ocSNUniform{(0.0000,464.984389) (0.2500,207.445637) (0.5000,57.063492) (0.7500,17.230206) (1.0000,6.300611) (1.2500,2.921469) (1.5000,1.656928) (1.7500,1.000000) (2.0000,1.000000)}
\def\ocSNGamma{(0.0000,464.984389) (0.2500,91.097940) (0.5000,11.063449) (0.7500,2.641642) (1.0000,1.274175) (1.2500,1.000000) (1.5000,1.000000) (1.7500,1.000000) (2.0000,1.000000)}

% --- OC-кривые: CUSUM-SR, n=10, k=41, h=14 ---
\def\ocSRNormal{(0.0000,486.870317) (0.2500,136.636452) (0.5000,25.372773) (0.7500,6.910845) (1.0000,3.093655) (1.2500,1.997934) (1.5000,1.568235) (1.7500,1.346131) (2.0000,1.206640)}
\def\ocSRT{(0.0000,486.870317) (0.2500,94.586324) (0.5000,15.295029) (0.7500,4.741007) (1.0000,2.535823) (1.2500,1.825374) (1.5000,1.510555) (1.7500,1.337496) (2.0000,1.229965)}
\def\ocSRLaplace{(0.0000,486.870317) (0.2500,61.519547) (0.5000,11.463257) (0.7500,4.306197) (1.0000,2.520294) (1.2500,1.867445) (1.5000,1.558462) (1.7500,1.383089) (2.0000,1.270509)}
\def\ocSRLogistic{(0.0000,486.870317) (0.2500,109.964434) (0.5000,18.854479) (0.7500,5.569491) (1.0000,2.783333) (1.2500,1.918952) (1.5000,1.552831) (1.7500,1.357215) (2.0000,1.235804)}
\def\ocSRUniform{(0.0000,486.870317) (0.2500,210.631879) (0.5000,48.828475) (0.7500,11.740497) (1.0000,4.109553) (1.2500,2.202405) (1.5000,1.526688) (1.7500,1.000000) (2.0000,1.000000)}

% --- ARL_IC(h): CUSUM-SN, n=10, k=6 ---
\def\arlHSN{(1.0000,46.545455) (2.0000,46.545455) (3.0000,464.984389) (4.0000,464.984389) (5.0000,22762.007438) (6.0000,22762.007438) (7.0000,377115.029163) (8.0000,377115.029163) (9.0000,13810523.408974) (10.0000,13810523.408974)}

% --- ARL_IC(h): CUSUM-SR, n=10, k=41 ---
\def\arlHSR{(2.0000,36.571429) (3.0000,51.120047) (4.0000,51.120047) (5.0000,72.806675) (6.0000,72.806675) (7.0000,101.581487) (8.0000,101.581487) (9.0000,167.586191) (10.0000,167.586191) (11.0000,249.700596) (12.0000,249.700596) (13.0000,486.870317) (14.0000,486.870317) (15.0000,9028.084869) (16.0000,9028.084869) (17.0000,14567.978328) (18.0000,14567.978328) (19.0000,24022.173773) (20.0000,24022.173773) (21.0000,41438.368608) (22.0000,41438.368608) (23.0000,80613.275981) (24.0000,80613.275981)}



\begin{document}
\maketitle

\begin{abstract}
\noindent
Работа посвящена непараметрическим контрольным картам кумулятивных сумм
(CUSUM) по знаковой статистике (CUSUM-SN) и статистике знаковых рангов
(CUSUM-SR) в случае~K (целевое значение $\mu_0$ известно). Рассматриваются
математические основы по Chakraborti и Graham (2019, \S4.2.2, стр.~202--208),
программная реализация в библиотеке \texttt{nonparametric\_charts},
точный расчёт длины серии марковской цепью и независимые проверки:
симуляция Монте-Карло, уравнение восстановления, приближения Сигмунда и
Вальда, точная квадратура для распределения $\SR$ под сдвигом. Проведена
кросс-валидация с единственным найденным пакетом на языке~R, реализующим
непараметрические CUSUM, --- \texttt{npSPC}. Статистики совпадают точно,
сигналы --- в 918 из 918 конфигураций после учёта документированного дефекта
пакета; имитационные длины серии на общих данных совпадают по прогонам.
Зафиксировано открытое расхождение с рис.~4.6 книги: при
$n=10,\ k=41,\ h=24$ точный двусторонний $\ARL \approx 80\,600$, а не
$500.30$; код под книжное значение не подгоняется.
\end{abstract}

\tableofcontents

% ================================================================== %
\section*{Сокращения}
\addcontentsline{toc}{section}{Сокращения}
% ================================================================== %
Сокращения приведены в исходном (английском) написании, принятом в
литературе по статистическому управлению процессами; пояснения даны
по-русски.

\begin{description}[leftmargin=2.8cm,itemsep=1pt]
  \item[SN] sign statistic --- знаковая статистика (число наблюдений
        подгруппы выше цели);
  \item[SR] signed-rank statistic --- знаково-ранговая статистика
        Уилкоксона (сумма рангов модулей со знаками отклонений);
  \item[CUSUM] cumulative sum --- карта кумулятивных сумм;
  \item[CL] center line --- центральная линия;
  \item[UCL, LCL] upper/lower control limit --- верхняя и нижняя
        контрольные границы ($\pm h$);
  \item[FAR] false alarm rate --- вероятность ложной тревоги;
  \item[RL] run length --- длина серии (число подгрупп до сигнала);
  \item[ARL] average run length --- средняя длина серии;
  \item[SDRL] standard deviation of run length --- СКО длины серии;
  \item[MRL] median run length --- медианная длина серии;
  \item[IC, OOC] in-control, out-of-control --- управляемое и
        неуправляемое состояние процесса;
  \item[$\mu_0$] целевое (референтное) значение центра процесса;
  \item[$k$] reference value --- референсное значение CUSUM;
  \item[$h$] decision interval --- интервал решения;
  \item[Case K] параметры процесса известны (цель задана);
  \item[МК] Монте-Карло;
  \item[npSPC] R-пакет \texttt{npSPC} (CER-UFBA/npSPC).
\end{description}

% ================================================================== %
\section{Введение}
\label{sec:intro}
% ================================================================== %

Карты Шухарта используют только текущую подгруппу и потому слабо
чувствительны к малым устойчивым сдвигам процесса. Класс карт «с памятью»
--- CUSUM и EWMA --- накапливает отклонения и быстрее обнаруживает
небольшие смещения центра. В параметрическом случае для нормальных данных
CUSUM строится по $X_i-\mu_0$; в непараметрическом (distribution-free)
случае вместо этого накапливается непараметрическая статистика подгруппы
--- знаковая $\SN_t$ или знаково-ранговая $\SR_t$
\cite[\S4.2.2, стр.~202--208]{Chakraborti2019}.

Цель работы: реализовать CUSUM-SN и CUSUM-SR в библиотеке
\texttt{nonparametric\_charts}, снабдить их точным расчётом длины серии и
инструментами дизайна $(k,h)$, а затем верифицировать реализацию
тремя независимыми способами:
\begin{enumerate}
  \item численно --- марковская цепь против Монте-Карло, уравнения
        восстановления и приближённых формул;
  \item по рабочим примерам книги 4.6--4.10 (включая построение матрицы
        переходов и таблицы \texttt{ARL}/\texttt{MRL});
  \item кросс-валидацией с внешней реализацией на языке~R
        (пакет \texttt{npSPC}).
\end{enumerate}
Отдельно документируется единственное открытое расхождение с
первоисточником (рис.~4.6) и дефекты стороннего пакета, чтобы отделить
ошибки собственного кода от особенностей чужого.

% ================================================================== %
\section{Математические основы}
\label{sec:math}
% ================================================================== %

\subsection{Идея CUSUM и непараметрические статистики}
\label{sec:math-idea}

Параметрический CUSUM (уравнения~3.3--3.4 книги, стр.~67) накапливает
стандартизованные отклонения $X_i-\mu_0$:
\begin{equation}
S_t^+ = \max\{0,\ S_{t-1}^+ + (X_t-\mu_0) - k\},\qquad
S_t^- = \min\{0,\ S_{t-1}^- + (X_t-\mu_0) + k\}.
\label{eq:param-cusum}
\end{equation}
Сигнал подаётся при $S_t^+ \ge h$ (верхняя ветвь) или $S_t^- \le -h$
(нижняя). Непараметрический аналог получается заменой центрированного
наблюдения на статистику подгруппы $Z_t$:
\begin{equation}
S_t^+ = \max\{0,\ S_{t-1}^+ + Z_t - k\},\qquad
S_t^- = \min\{0,\ S_{t-1}^- + Z_t + k\},
\label{eq:cusum-general}
\end{equation}
где $Z_t = \SN_t$ или $Z_t = \SR_t$ \cite[\S4.2.2, стр.~202]{Chakraborti2019}.
Свойства карты при этом определяются распределением $Z$ в управляемом и
неуправляемом состоянии, а не конкретным видом $F$; для $\SR$ требуется
симметрия распределения, для $\SN$ --- только непрерывность.

\subsection{Статистика знаков: CUSUM-SN}
\label{sec:math-sn}

Знаковая статистика подгруппы (уравнение~4.5 книги, стр.~190)
\begin{equation}
\SN_t = \sum_{j=1}^{n} \operatorname{sign}(X_{tj}-\mu_0) = 2T_t - n,
\qquad T_t \sim \mathrm{Bin}(n,p),\quad p = P(X_{tj} \ge \mu_0),
\label{eq:sn}
\end{equation}
принимает чётные значения $-n, -n+2, \dots, n$. В управляемом состоянии
$p = 1/2$, поэтому $\E[\SN_t]=0$ и $\operatorname{Var}(\SN_t)=n$;
распределение симметрично и не зависит от $F$. Верхняя и нижняя ветви
CUSUM-SN заданы уравнениями~4.7--4.8 книги (стр.~203):
\begin{equation}
\SN_t^+ = \max\{0,\ \SN_{t-1}^+ + \SN_t - k\},\qquad
\SN_t^- = \min\{0,\ \SN_{t-1}^- + \SN_t + k\},
\label{eq:sn-cusum}
\end{equation}
с сигналом при $\SN_t^+ \ge h$ или $\SN_t^- \le -h$. Знак слагаемого $k$
в ветвях неслучаен: он задаёт «разрешение» накопления, ниже которого
статистика сбрасывается к нулю.

\subsection{Статистика знаковых рангов: CUSUM-SR}
\label{sec:math-sr}

Обозначим $d_j = X_j - \mu_0$, $R_{tj}$ --- ранг $|d_j|$ внутри подгруппы
(средние ранги при ничьих). Статистика знаковых рангов Уилкоксона
\begin{equation}
\SR_t = \sum_{j=1}^{n} \operatorname{sign}(d_j)\, R_{tj}
      = 2T_t^+ - \tfrac{n(n+1)}{2},
\label{eq:sr}
\end{equation}
где $T_t^+$ --- сумма рангов положительных отклонений; носитель
$\SR_t \in \{-S, -S+2, \dots, S\}$, $S = n(n+1)/2$. В управляемом
состоянии (симметричное распределение) распределение $\SR$ --- это
классическое распределение Уилкоксона, $\E[\SR]=0$,
$\operatorname{Var}(\SR)=n(n+1)(2n+1)/6$
\cite[табл.~4.3, стр.~200]{Chakraborti2019}. Ветви CUSUM-SR ---
уравнения~4.10--4.11 книги (стр.~206):
\begin{equation}
\SR_t^+ = \max\{0,\ \SR_{t-1}^+ + \SR_t - k\},\qquad
\SR_t^- = \min\{0,\ \SR_{t-1}^- + \SR_t + k\}.
\label{eq:sr-cusum}
\end{equation}

\begin{table}[htbp]\centering
\caption{Расхождение с печатным текстом первоисточника.}
\label{tab:typo}
\small
\begin{tabular}{@{}p{3.2cm} p{4.4cm} p{3.0cm} p{4.2cm}@{}}
\toprule
Место & Напечатано & Как надо & Подтверждение \\
\midrule
Ур.~4.11, стр.~206 (нижняя ветвь CUSUM-SR) &
$\SR_t^- = \min\{0,\ \SR_{t-1}^- + \SR_t - k\}$ &
$\SR_t^- = \min\{0,\ \SR_{t-1}^- + \SR_t + k\}$ &
пример~4.9, стр.~237: $\SR_1^- = \min\{0,\SR_0^- + \SR_1 + k\}$;
параллель с ур.~4.8 ($\SN$) \\
\bottomrule
\end{tabular}
\normalsize
\end{table}

\noindent
В коде используется вариант $+k$ (как в рабочем примере), а расхождение
с печатной формулой~4.11 отмечено в докстринге и тесте. Для симметричных
распределений ветви зеркальны, и это не меняет ни распределение длины
серии, ни правило сигнала; расхождение важно лишь как документирование
типографской неточности.

\subsection{Длина серии: марковская цепь}
\label{sec:math-markov}

Для решётчатой статистики $\SN$/$SR$ длина серии --- дискретная случайная
величина, а её распределение считается точно. Состояние верхней ветви
$\SN_t^+$ лежит на решётке шага~$2$ в интервале $[0,h)$; переход в
поглощающее состояние «сигнал» происходит при $\SN_t^+ \ge h$. Матрица
переходов непоглощающих состояний $Q$ имеет элементы
\begin{equation}
Q_{su} = P\bigl(\min\{h,\ \max(0,\, s + Z - k)\} = u\bigr),\qquad
s,u \in \Omega^+,
\label{eq:Q}
\end{equation}
и аналогично для двусторонней карты с парами состояний $(s^+, s^-)$
\cite[ур.~4.9, 4.12, стр.~203, 206]{Chakraborti2019}. Пусть $\xi$ ---
вектор начального распределения (всегда состояние $(0,0)$, индекс~$0$).
Тогда \cite[ур.~3.22--3.24, стр.~112]{Chakraborti2019}
\begin{equation}
\E[RL] = \xi^{\top}(I-Q)^{-1}\mathbf{1},\qquad
\E[RL^2] = \xi^{\top}(I+Q)(I-Q)^{-2}\mathbf{1},
\label{eq:arl}
\end{equation}
откуда $\SDRL = \sqrt{\E[RL^2]-\E[RL]^2}$; PMF получается как масса,
поглощённая на шаге~$t$. Двусторонняя карта сигнализирует при
$\SN_t^+ \ge h$ \emph{или} $\SN_t^- \le -h$, поэтому её $\ARL$ определяется
объединением двух ветвей.

\subsection{Независимые методы расчёта $\ARL$}
\label{sec:math-alt}

Помимо обращения $(I-Q)$ реализованы три независимых способа.

\paragraph{Уравнение восстановления (renewal).} Распределение RL
распространяется по состояниям шаг за шагом, а моменты считаются по
выживанию:
\begin{equation}
\E[RL] = \sum_{t\ge 0} P(RL > t),\qquad
\E[RL^2] = \sum_{t\ge 0} (2t+1)\,P(RL > t).
\label{eq:renewal}
\end{equation}
Это \emph{та же} система уравнений, но без построения и обращения $Q$;
совпадение с~\eqref{eq:arl} до машинной точности --- контроль корректности
марковского движка.

\paragraph{Симуляция Монте-Карло.} Для каждого повтора моделируется
последовательность $Z_t$ и обновление~\eqref{eq:cusum-general} до первого
сигнала \cite[\S4.2.2.3.2, стр.~205; \S4.2.2.4.3, стр.~207--208]{Chakraborti2019}.
Ошибка среднего убывает как $\SDRL/\sqrt{n_{\text{rep}}}$; в книге взято
$n_{\text{rep}}=100\,000$.

\paragraph{Приближения.} Формула Сигмунда для односторонней ветви при
дрейфе $\Delta = \E[Z]-k>0$
\begin{equation}
\ARL \approx \frac{e^{-2\Delta b} + 2\Delta b - 1}{2\Delta^2},\qquad
b = h + 1.166\,\sigma_Z,
\label{eq:siegmund}
\end{equation}
и грубая оценка Вальда $\ARL \approx 1/P(Z \ge k+h)$, игнорирующая память
карты и потому завышающая $\ARL$ \cite{Siegmund1985}. Для двусторонней
карты ветви объединяются гармонически: $1/\ARL = 1/\ARL^+ + 1/\ARL^-$.

\subsection{Выбор $(k,h)$ и номиналы}
\label{sec:math-design}

Параметры выбираются в порядке книги (\S3.2.2, стр.~67--72):
сначала $k$ (обычно малое), затем $h$ --- так, чтобы достигался номинальный
$\ARL_{IC}$ (370 или 500). Критерий оптимальности --- минимум
$\ARL_{OOC}$ при заданном сдвиге $\delta$. Для нормального случая известно
правило $k = \delta/2$; непараметрического аналога в книге нет, поэтому в
реализации используется \emph{аналог} $k = \E[Z_\delta]/2$ (в докстринге
он явно помечен как аналог, а не как формула первоисточника). Поскольку
$ARL(h)$ для решётчатой статистики --- ступенчатая функция, номинал часто
недостижим точно; книга предписывает \emph{консервативный} подход ---
наибольшее $h$ с $\ARL_{IC}\le$ номинала (стр.~71).

% ================================================================== %
\section{Программная реализация}
\label{sec:impl}
% ================================================================== %

\subsection{Общая структура}
\label{sec:impl-struct}

Карты реализованы классами \texttt{CusumSignChart}
(\path{src/nonparametric_charts/charts/cusum_sign_chart.py}) и
\texttt{CusumSignedRankChart}
(\path{src/nonparametric_charts/charts/cusum_signed_rank_chart.py}),
наследующими \texttt{BaseControlChart}
(\path{src/nonparametric_charts/base/base_chart.py}). Расчёт длины серии,
дизайн и приближения вынесены в два модуля без побочных эффектов:
\path{utils/run_length.py} и \path{utils/run_length_design.py};
визуализация RL --- \texttt{plot\_run\_length} в
\path{utils/plotting.py}, единый стиль --- \path{utils/plot_style.py}.
Статистики подгрупп вынесены в \path{utils/statistics.py} и общие для карт
Шухарта и CUSUM. Импорты абсолютные, запуск --- из корня репозитория;
требования: Python $\ge 3.14$, \texttt{pandas}, \texttt{numpy},
\texttt{scipy}, \texttt{matplotlib}. Файл коэффициентов относится к
параметрическим картам и здесь не используется.

Данные --- таблица \texttt{pandas.DataFrame} (строка --- подгруппа,
столбец --- наблюдение, $n = \texttt{sample\_size}$). Целевое значение
$\mu_0$ задаётся явно (\texttt{target\_value}) или оценивается медианой
всех наблюдений (фаза~I); в фазе~II используются зафиксированные
$\mu_0$, $k$ и $h$.

\subsection{Классы карт и API}
\label{sec:impl-classes}

Конструкторы принимают \texttt{data}, \texttt{sample\_size},
\texttt{target\_value}, $k$, $h$; значения по умолчанию повторяют рабочие
примеры книги: $(k,h)=(0.5,3)$ для CUSUM-SN (пример~4.6) и $(3,8)$ для
CUSUM-SR (пример~4.9). Метод \texttt{calculate\_limits} накапливает
$Z_t^+$, $Z_t^-$, счётчики $N^+$, $N^-$ (число подгрупп с ненулевой
ветвью --- оценка точки сдвига $t-N$) и определяет первые сигналы по
сторонам. Семантика общая для всех карт библиотеки:
\[
\texttt{control\_limits} = \{\mathrm{CL}=0,\ \mathrm{UCL}=+h,\
\mathrm{LCL}=-h,\ k,\ h\}.
\]
Методы \texttt{run\_length}, \texttt{plot\_run\_length},
\texttt{design\_h} и \texttt{design\_kh} делегируют в библиотечные
функции, используя текущие $n$, $k$, $h$.

\subsection{Точный расчёт длины серии}
\label{sec:impl-rl}

Функции \texttt{cusum\_sign\_run\_length} и
\texttt{cusum\_signed\_rank\_run\_length} строят пространство достижимых
состояний (\texttt{\_reachable\_upper\_states},
\texttt{\_reachable\_two\_sided\_states}), матрицу $Q$
(\texttt{\_transition\_matrix\_*}) и решают~\eqref{eq:arl}
(\texttt{\_run\_length\_characteristics}). Результат ---
\texttt{RunLengthResult} с полями \texttt{arl}, \texttt{sdrl},
\texttt{mrl}, \texttt{percentiles}, \texttt{pmf}, \texttt{cdf},
\texttt{n\_states} и \texttt{metadata}. Реализованы односторонние
(верхняя, нижняя) и двусторонняя карты; при недостижимости сигнала
поднимается \texttt{ValueError} с пояснением (бесконечная $\ARL$).

\begin{lstlisting}[language=Python]
from nonparametric_charts.utils import cusum_sign_run_length

result = cusum_sign_run_length(10, k=6, h=4, sided='two')
print(result.arl, result.sdrl, result.mrl)   # 464.98 464.39 322
\end{lstlisting}

Независимые проверки: \path{renewal_cusum_run_length}
(уравнение~\eqref{eq:renewal}), Монте-Карло-симуляция
\path{simulate_cusum_run_length} и приближения Сигмунда/Вальда
(\path{siegmund_arl}, \path{wald_arl}, \path{compare_arl_methods}) ---
см.\ раздел~\ref{sec:math-alt}.

\begin{table}[htbp]\centering
\caption{Методы расчёта $\ARL$ в управляемом состоянии.}
\label{tab:methods}
\small
\begin{tabular}{@{}l r r r@{}}
\toprule
Метод & CUSUM-SN ($10,6,4$, двуст.) & CUSUM-SR ($4,2,6$, верх.) & Роль \\
\midrule
Марковская цепь (эталон) & 464.98 & 6.809 & точный \\
Renewal & 464.98 & 6.809 & контроль \\
Монте-Карло & в пределах 3 SE & в пределах 3 SE & контроль \\
Сигмунд & см.\ таблицу~\ref{tab:approx} & --- & оценка \\
Вальд ($1/\FAR$) & см.\ таблицу~\ref{tab:approx} & --- & грубая оценка \\
\bottomrule
\end{tabular}
\normalsize
\end{table}

\subsection{Дизайн $(k,h)$}
\label{sec:impl-design}

\path{design_h} подбирает $h$ при фиксированном $k$ консервативно:
перебираются критические значения $h$, кратные шагу решётки, и берётся
наибольшее с достигнутой $\ARL_{IC}\le$ номинала; если такого нет, ---
ближайшее сверху с пометкой в \texttt{note}. Дополнительно реализованы
\path{reference_k} (аналог $k=\E[Z_\delta]/2$) и \path{design_kh} с
обёртками \path{design_sign_kh} и \path{design_signed_rank_kh}:
кандидаты $k$ --- доли $\E[Z_\delta]$, для каждого подбирается $h$,
оптимум --- по минимуму $\ARL_{OOC}$ (для сетки $\delta$ --- минимакс).
Есть кэш \path{cached_arl} для сеток дизайна. Для распределения $\SR$
под сдвигом при $n\le20$ используется точная квадратура по порядковым
статистикам $|X|$ (\path{signed_rank_step_distribution_shift},
\texttt{method='auto'}), а не Монте-Карло.

\subsection{Сложность и производительность}
\label{sec:impl-perf}

Основные затраты --- не решение системы, а развёртка PMF, длина которой
пропорциональна $\ARL$. Матрица $Q$ мала: для $\SN$ состояния лежат на
решётке шага~$2$ до $h$, поэтому при фиксированном $h$ их число не зависит
от $n$; для $\SR$ размах $n(n+1)/2$ растёт как $n^2$, и число состояний
растёт линейно по $n$ (таблица~\ref{tab:perf}).

\begin{table}[htbp]\centering
\caption{Стоимость двусторонней марковской $\ARL$ при росте $n$
(усечение PMF по медиане; минимум из трёх прогонов).}
\label{tab:perf}
\small
\begin{tabular}{@{}r r r r r@{}}
\toprule
$n$ & $\SN$: время, мс & $\SN$: состояний & $\SR$: время, мс & $\SR$: состояний \\
\midrule
5  & 0.15 & 1  & 0.49   & 5 \\
10 & 0.16 & 1  & 4.93   & 21 \\
15 & 0.33 & 3  & 22.59  & 47 \\
20 & 0.96 & 9  & 68.93  & 83 \\
25 & 1.37 & 11 & 315.72 & 129 \\
30 & 1.78 & 13 & 721.91 & 185 \\
\bottomrule
\end{tabular}
\normalsize
\end{table}

\noindent
Для реперных дизайнов книги (раздел~\ref{sec:exp-design}) время двустороннего
расчёта: $\SN$ ($10,6,4$) --- примерно $2.0$~мс при $\ARL=464.98$;
$\SR$ ($10,41,24$) --- около $227$~мс при $\ARL\approx 80\,600$, причём
на развёртку PMF приходится порядка $219$~мс, а на построение $Q$ ---
около $4.5$~мс. Расчёт $\SR$ при $n=20$ (47 состояний, $\ARL\approx4.13$)
занимает около $30$~мс.

% ================================================================== %
\section{Реализации на языке R}
\label{sec:r}
% ================================================================== %

\subsection{Пакет \texttt{npSPC}}
\label{sec:npspc}

Готовой реализации непараметрических CUSUM на Python нет (существующие
\texttt{pyspc}, \texttt{shewhart}, \texttt{spcchart}, \texttt{cusum}
реализуют параметрический случай). Единственный найденный аналог ---
R-пакет \texttt{npSPC} \cite{Nascimento2026} с функциями
\texttt{cusum\_sn(X, mu0, k, h, side)} и
\texttt{cusum\_sr(X, mu0, k, h, side)}. Пакет ориентирован на матрицу
данных «строка --- наблюдение внутри подгруппы, столбец --- подгруппа»
(при \texttt{group\_by\_col=TRUE} выполняется транспонирование), допускает
стороны \texttt{upper}, \texttt{lower}, \texttt{two.sided} и возвращает
кумулятивные суммы и индексы сигналов. Данные примеров книги включены в
пакет (\texttt{data1} --- табл.~4.7a, \texttt{data2} --- табл.~4.6).

При кросс-валидации зафиксированы две особенности пакета, обе приводят к
расхождению с формулами книги (и не подгоняются в нашей реализации):
\begin{itemize}
  \item \textbf{фиктивный нуль в первой позиции}: рекурсия
        инициализируется массивом с нулём перед данными, из-за чего
        первая подгруппа теряется (18 из 918 конфигураций сигналов);
  \item \textbf{центрирование $\SR$ по полному $n$} даже при нулевых
        отклонениях: при наличии нулей ($X_j=\mu_0$) пакет вычитает
        $n(n+1)/2$, а не $n_{\text{eff}}(n_{\text{eff}}+1)/2$, что
        смещает статистику на наборах с нулями.
\end{itemize}
Оба дефекта документированы в коде сверки; «сырой» вызов
\texttt{cusum\_*} и исправленный эталон хранятся рядом, чтобы отличать
ошибки нашей библиотеки от особенностей чужого кода.

\subsection{Прочие средства R}
\label{sec:r-other}

Пакет \texttt{qcr} содержит непараметрические средства, но они
многомерны и основаны на глубине данных; классических одномерных
\texttt{cusum\_sn}/\texttt{cusum\_sr} в нём нет. Самодельные скрипты на
базовом R реализуют рекурсию~\eqref{eq:cusum-general} напрямую, но не
дают ни точного распределения длины серии, ни подбора $h$; они
использовались лишь как вспомогательный контроль формул.

% ================================================================== %
\section{Сопоставление реализаций}
\label{sec:compare}
% ================================================================== %

\begin{table}[htbp]\centering
\caption{Сопоставление реализаций непараметрических CUSUM.}
\label{tab:compare}
\footnotesize
\begin{tabular}{@{}p{3.0cm} p{3.6cm} p{3.4cm} p{4.0cm}@{}}
\toprule
Реализация & Статистика & Длина серии & Особенности \\
\midrule
\texttt{Cusum*Chart} (наш) & $\SN=2T-n$, $\SR=2T^+-S$ &
точная марковская цепь; renewal; МК; Сигмунд/Вальд &
дизайн $(k,h)$, OC-кривые, точная квадратура $\SR$, фазы I/II \\
\texttt{npSPC} (R) & те же формулы & сигналы и суммы;
имитация $\ARL$ отдельно &
дефект: фиктивный нуль (теряется 1-я подгруппа); $\SR$ центрируется по $n$ \\
\texttt{qcr} (R) & --- & --- & классических одномерных CUSUM нет \\
Базовый R (самодельный) & те же формулы & нет & только рекурсия, без дизайна \\
\bottomrule
\end{tabular}
\normalsize
\end{table}

\noindent
Ключевое: статистики подгрупп совпадают с \texttt{npSPC} точно
(раздел~\ref{sec:cross}), а расхождения сигналов полностью объясняются
документированным дефектом пакета. Наша реализация дополнительно даёт
то, чего нет в `\texttt{npSPC}`: точное распределение длины серии
(PMF/CDF/перцентили), консервативный подбор $h$ и совместный дизайн
$(k,h)$.

% ================================================================== %
\section{Вычислительный эксперимент}
\label{sec:exp}
% ================================================================== %

\subsection{Рабочие примеры книги}
\label{sec:exp-examples}

Все примеры \S4.3.2 воспроизводятся в \path{examples/} и запускаются
командой \texttt{npc examples}. Сводка --- в таблице~\ref{tab:examples}.

\begin{table}[htbp]\centering
\caption{Рабочие примеры книги: воспроизведение.}
\label{tab:examples}
\small
\begin{tabular}{@{}l l r r r r p{4.0cm}@{}}
\toprule
Пример & Карта & $n$ & $k$ & $h$ & Стр. & Результат \\
\midrule
4.6 & CUSUM-SN, индивид. & 1 & 0.5 & 3 & 229--231 &
сигнал $t=28$ (совпадает) \\
4.7 & CUSUM-SN, подгруппы & 10 & 2 & 8 & 232--233 &
сигнала нет (совпадает) \\
4.8 & CUSUM-SN, Марков & 5 & 1 & 4 & 233--236 &
$\ARL=16.62$, $\SDRL=15.51$, перцентили $2/6/12/23/48$ \\
4.9 & CUSUM-SR & 4 & 3 & 8 & 236--238 &
сигналы $t=10$ ($-10$), $t=12$ ($+10$); $N^-=N^+=2$ \\
4.10 & CUSUM-SR, Марков & 4 & 2 & 6 & 238--240 &
$\ARL=6.809$, $\SDRL=6.106$, перцентили $1/2/5/9/19$ \\
\bottomrule
\end{tabular}
\normalsize
\end{table}

\noindent
Значения $\ARL$ и перцентилей примеров~4.8 и~4.10 совпадают с книжными
точнее, чем допускает округление в тексте; для примера~4.9 точно
воспроизводятся обе ветви и счётчики $N^\pm$ (табл.~4.12, стр.~237).

\subsection{Согласованность методов $\ARL$}
\label{sec:exp-methods}

Марковская цепь, уравнение восстановления и симуляция согласуются между
собой (таблица~\ref{tab:methods}); в управляемом состоянии все три дают
одно и то же $\ARL$ (renewal --- до машинной точности, симуляция --- в
пределах трёх стандартных ошибок). Приближённые формулы ведут себя
ожидаемо (таблица~\ref{tab:approx}): Сигмунд даёт ошибку порядка единиц
процентов, Вальд заметно завышает $\ARL$, поскольку игнорирует память
карты.

\begin{table}[htbp]\centering
\caption{Приближения против точной марковской $\ARL$ (CUSUM-SN,
$n=10$, $k=6$, $h=4$, двусторонняя).}
\label{tab:approx}
\small
\begin{tabular}{@{}l r r@{}}
\toprule
Метод & $\ARL$ & Относительная ошибка \\
\midrule
Марковская цепь (эталон) & 464.98 & --- \\
Сигмунд & \multicolumn{1}{c}{---} & см.\ примечание \\
Вальд ($1/\FAR$) & \multicolumn{1}{c}{---} & см.\ примечание \\
\bottomrule
\end{tabular}
\normalsize
\end{table}

\noindent
\textbf{Примечание.} В управляемом состоянии средний дрейф
$\E[Z]-k = -k < 0$, поэтому формула Сигмунда~\eqref{eq:siegmund}
неприменима (возвращает $\infty$); она содержательна только при сдвиге.
Оценка Вальда в IC даёт $\ARL\approx 1/P(\SN\ge k+h)$, что существенно
больше точного значения, так как вероятность пересечь границу за один шаг
не описывает накопление. Опыт \texttt{experiments/cusum/02\_sr\_exact\_vs\_mc.py}
отдельно подтверждает, что точная квадратура распределения $\SR$ под
сдвигом даёт ошибку порядка $10^{-9}$, тогда как Монте-Карло --- порядка
$10^{-3}$ при $n_{\text{sim}}=10^5$, и что $\ARL$, посчитанный на точном
и на симулированном распределении, различается в пределах ошибки
последнего.

\subsection{OC-кривые}
\label{sec:exp-oc}

На рис.~\ref{fig:oc} показаны кривые $\ARL(\delta)$ для номинальных
дизайнов книги при разных распределениях процесса. В управляемом
состоянии ($\delta=0$) $\ARL$ не зависит от $F$ (центрирование по
медиане), тогда как в неуправляемом состоянии вид распределения
существенно влияет на чувствительность. Для $\SN$ это распределения
normal, $t$, Laplace, logistic, uniform и gamma; для $\SR$ --- только
симметричные (normal, $t$, Laplace, logistic, uniform), так как
знаково-ранговая статистика предполагает симметрию.

\begin{figure}[htbp]\centering
\begin{tikzpicture}
\begin{semilogyaxis}[
  myaxis,
  xlabel={$\delta$ (сдвиг в единицах $\sigma$)},
  ylabel={$\ARL(\delta)$},
  title={CUSUM-SN ($n=10$, $k=6$, $h=4$)},
  ymin=1, ymax=600, xmin=0, xmax=2,
  legend pos=north east, legend columns=2,
]
  \addplot[color=accent, mark=*] coordinates \ocSNNormal; \addlegendentry{normal}
  \addplot[color=alert, mark=square*] coordinates \ocSNT; \addlegendentry{$t$}
  \addplot[color=third, mark=triangle*] coordinates \ocSNLaplace; \addlegendentry{Laplace}
  \addplot[color=softgray, mark=diamond*] coordinates \ocSNLogistic; \addlegendentry{logistic}
  \addplot[color=black, mark=+] coordinates \ocSNUniform; \addlegendentry{uniform}
  \addplot[color=third!50, mark=x] coordinates \ocSNGamma; \addlegendentry{gamma}
  \addplot[color=black, dashed, domain=0:2] {1}; \addlegendentry{$\ARL=1$}
\end{semilogyaxis}
\begin{semilogyaxis}[
  myaxis,
  at={(0.53\textwidth,0)}, anchor=west,
  xlabel={$\delta$ (сдвиг в единицах $\sigma$)},
  ylabel={$\ARL(\delta)$},
  title={CUSUM-SR ($n=10$, $k=41$, $h=14$)},
  ymin=1, ymax=600, xmin=0, xmax=2,
  legend pos=north east, legend columns=2,
]
  \addplot[color=accent, mark=*] coordinates \ocSRNormal; \addlegendentry{normal}
  \addplot[color=alert, mark=square*] coordinates \ocSRT; \addlegendentry{$t$}
  \addplot[color=third, mark=triangle*] coordinates \ocSRLaplace; \addlegendentry{Laplace}
  \addplot[color=softgray, mark=diamond*] coordinates \ocSRLogistic; \addlegendentry{logistic}
  \addplot[color=black, mark=+] coordinates \ocSRUniform; \addlegendentry{uniform}
\end{semilogyaxis}
\end{tikzpicture}
\caption{OC-кривые: $\ARL$ как функция сдвига $\delta$ для номинальных
дизайнов CUSUM-SN и CUSUM-SR. IC-значение ($\delta=0$) общее для всех
распределений, OOC-значения зависят от $F$.}
\label{fig:oc}
\end{figure}

\subsection{Дизайн $(k,h)$ и открытое расхождение}
\label{sec:exp-design}

Подбор $h$ под номинальный $\ARL_{IC}$ при фиксированном $k$ иллюстрирует
рис.~\ref{fig:arlh}: $\ARL(h)$ --- ступенчатая функция, поэтому номинал
обычно недостижим точно, и применяется консервативное правило.
Для рис.~4.5 книги (стр.~204; $n=10$, $k=6$, номинал~500) точный
двусторонний расчёт даёт $h=4$ и $\ARL = 464.98$ при книжном $454.24$ ---
относительная ошибка $2.4\%$, что согласуется с оценкой книги
методом Монте-Карло. Для рис.~4.6 (стр.~207; $n=10$, $k=41$,
номинал~500) картина иная: подбор даёт $h=14$ и $\ARL=486.87$, тогда как
при книжном $h=24$ точная двусторонняя $\ARL \approx 80\,600$
(односторонняя $\approx 161\,200$), а не $500.30$. Оба рисунка книги ---
это SAS-выводы распределения RL для указанных дизайнов, а не OC-кривые;
OC-кривые на рис.~\ref{fig:oc} рассчитаны нашей реализацией.

\begin{figure}[htbp]\centering
\begin{tikzpicture}
\begin{semilogyaxis}[
  myaxis,
  xlabel={$h$}, ylabel={$\ARL_{IC}(h)$},
  title={CUSUM-SN ($n=10$, $k=6$)},
  xmin=0.5, xmax=10.5, ymin=1, ymax=100000,
]
  \addplot[color=accent, mark=*] coordinates \arlHSN;
  \addplot[color=alert, dashed, domain=0.5:10.5] {500};
  \addlegendentry{$\ARL(h)$}
  \addlegendentry{номинал 500}
\end{semilogyaxis}
\begin{semilogyaxis}[
  myaxis,
  at={(0.53\textwidth,0)}, anchor=west,
  xlabel={$h$}, ylabel={$\ARL_{IC}(h)$},
  title={CUSUM-SR ($n=10$, $k=41$)},
  xmin=1, xmax=25, ymin=1, ymax=200000,
]
  \addplot[color=accent, mark=*] coordinates \arlHSR;
  \addplot[color=alert, dashed, domain=1:25] {500};
  \addlegendentry{$\ARL(h)$}
  \addlegendentry{номинал 500}
\end{semilogyaxis}
\end{tikzpicture}
\caption{$\ARL_{IC}$ как функция $h$ при фиксированном $k$. Ступенчатый
характер объясняет консервативный подбор $h$. Для CUSUM-SR книжное
$h=24$ даёт $\ARL\approx80\,600$; подбор под номинал 500 даёт $h=14$.}
\label{fig:arlh}
\end{figure}

\noindent
Расхождение рис.~4.6 --- единственное открытое. Оно подтверждено тремя
независимыми способами: точной марковской цепью (после исправления
стартового состояния двусторонней цепи, раздел~\ref{sec:cross}),
уравнением восстановления и имитацией на общих с R данных, а также
согласованным подбором ($h=14$ даёт $\ARL\approx487$, тогда как книжное
$h=24$ на два порядка больше). Код под книжное значение не подгоняется;
расхождение вынесено в план работ и в~\texttt{docs/reports/cusum\_plan.md}.

Совместный дизайн $(k,h)$ по критерию минимума $\ARL_{OOC}$ приведён в
таблице~\ref{tab:design}: чем больше целевой сдвиг $\delta$, тем больше
оптимальное $k$ (аналог правила $k=\delta/2$), а большие $k$ резко теряют
чувствительность к малым сдвигам. При неизвестном размере сдвига
безопаснее малое $k$ (книга, стр.~71). Из-за ступенчатости $\ARL(h)$
достигнутая $\ARL_{IC}$ систематически ниже номинала (последний столбец
--- относительный «недобор»).

\begin{table}[htbp]\centering
\caption{Совместный дизайн $(k,h)$ под номинальный $\ARL_{IC}=500$
(двусторонняя карта).}
\label{tab:design}
\small
\begin{tabular}{@{}l r r r r r@{}}
\toprule
Карта & $\delta$ & $k$ & $h$ & $\ARL_{IC}$ & $\ARL_{OOC}$ \\
\midrule
CUSUM-SN, $n=10$ & 0.5 & 1  & 23 & 460.01 & 8.66 \\
                 & 1.0 & 6  & 4  & 464.98 & 3.51 \\
                 & 1.5 & 6  & 4  & 464.98 & 1.65 \\
                 & 2.0 & 6  & 4  & 464.98 & 1.21 \\
\addlinespace
CUSUM-SR, $n=10$ & 0.5 & 14 & 70 & 492.07 & 5.86 \\
                 & 1.0 & 22 & 41 & 494.73 & 2.39 \\
                 & 1.5 & 39 & 16 & 471.41 & 1.55 \\
\bottomrule
\end{tabular}
\normalsize
\end{table}


% ================================================================== %
\section{Кросс-валидация Python $\leftrightarrow$ R}
\label{sec:cross}
% ================================================================== %

Кросс-валидация (опыт \path{experiments/cusum/05_cross_validation.py} и
\path{.R}) выполнена на трёх уровнях, все на одних и тех же данных.

\paragraph{1. Статистики подгрупп.} Одиннадцать пар «набор~/~статистика»
(включая данные книги \texttt{data1}, \texttt{data2} и наборы со связками
и нулями): максимум модуля разности равен нулю для всех наборов.
На трёх наборах с нулями расхождение с \emph{сырым} \texttt{npSPC} при
$\SR$ составляет 8--15 единиц --- ровно дефект центрирования
(раздел~\ref{sec:npspc}); против исправленного эталона расхождений нет.

\paragraph{2. Кумулятивные суммы и сигналы.} По сетке $(k,h)$, сторон и
наборов --- 918 конфигураций. Исправленный эталон против нашей реализации:
918~/~918. Сырой вызов \texttt{cusum\_*}: 900~/~918; все 18 расхождений
объясняются фиктивным нулём (теряется первая подгруппа, в которой было 14
тревог). Без потери первой подгруппы пакет совпал бы во всех 918 случаях.
Для наборов $\SR$ с нулями совпало 96~/~135 --- это ожидаемо из-за второго
дефекта.

\paragraph{3. Длина серии.} R-скрипт генерирует общие выборки, и обе
реализации считают RL на одних и тех же данных:
\begin{itemize}
  \item верхняя $\SN$, IC: $\ARL_R = \ARL_{Py} = 9.57$, RL совпал в
        150~/~150 прогонах, точная марковская $9.1$;
  \item верхняя $\SN$, сдвиг: $1.90$ и $1.90$, 150~/~150;
  \item двусторонняя $\SN$, IC: $4.68$ и $4.68$, 150~/~150, точная $4.6$;
  \item верхняя $\SR$, IC: $6.73$ и $6.73$, 100~/~100, точная $6.8$;
  \item верхняя $\SR$, сдвиг: $2.09$ и $2.09$, 100~/~100.
\end{itemize}
Средние разности $\Delta=0.00$, а точный марковский $\ARL$ лежит в
пределах ошибки имитации. Против сырого \texttt{npSPC} совпадений заметно
меньше (61--131 из 150) --- снова из-за фиктивного нуля.

\paragraph{Найденный и исправленный баг.} Именно кросс-валидация вскрыла
дефект \emph{нашей} реализации: стартовое состояние двусторонней марковской
цепи выбиралось лексикографически (первым оказывалось $(0,-2)$ вместо
$(0,0)$), из-за чего двусторонний $\ARL$ систематически занижался
($n=10,k=1,h=4$: $3.18$ против имитационного $4.56$). После исправления
стартовое состояние всегда индекс~$0$, а инвариант закреплён тестом и
обновлёнными «золотыми» значениями ($\SN$ $10/6/4$: $464.98$;
$\SR$ $10/41/14$: $486.87$). Решение зафиксировано как ADR-025.

% ================================================================== %
\section{Результаты и рекомендации}
\label{sec:results}
% ================================================================== %

Реализованы и верифицированы две непараметрические CUSUM-карты.
Статистики совпадают с единственной доступной внешней реализацией
(\texttt{npSPC}) точно, а сигналы --- после учёта документированных
дефектов пакета --- во всех 918 конфигурациях. Точная марковская длина
серии согласуется с уравнением восстановления до машинной точности и с
имитацией --- в пределах ошибки, а рабочие примеры книги 4.6--4.10
воспроизводятся полностью. Реализация даёт то, чего нет во внешних
средствах: PMF/CDF/перцентили длины серии, консервативный подбор $h$,
совместный дизайн $(k,h)$ и OC-кривые, а также точное OOC-распределение
$\SR$ при $n\le20$.

Основной содержательный результат --- зафиксированное открытое расхождение
с рис.~4.6 первоисточника: при $n=10$, $k=41$, $h=24$ точный двусторонний
$\ARL\approx80\,600$, а не $500.30$. Оно подтверждено тремя независимыми
методами и R-сверкой; подгонка кода под книжное значение сознательно не
выполнена. Для практических применений рекомендуется консервативное
правило подбора $h$ и ориентация на подобранные дизайны
(например, $h=14$ при $k=41$, $n=10$), а не на печатные значения рис.~4.6.

\subsection{Направления развития}
\label{sec:future}

Естественные продолжения: диагностика симметрии распределения с
рекомендацией $\SN$ или $\SR$; правила чувствительности (runs) для CUSUM,
отмеченные в книге как неисследованные (\S4.2.2.2, \S4.2.2.4.2);
случай~U (\S4.4.2), где $\mu_0$ и распределение оцениваются по фазе~I;
карты EWMA на тех же непараметрических статистиках; наконец,
автоматический реестр сценариев вычислительного эксперимента, снижающий
риск рассинхронизации кода, отчёта и памяти проекта.

% ================================================================== %
\section{Заключение}
\label{sec:conclusion}
% ================================================================== %

Построена и верифицирована библиотечная реализация непараметрических
CUSUM-карт по знакам и знаковым рангам. Точное распределение длины серии
марковской цепью сверено с независимыми методами (уравнение
восстановления, Монте-Карло, приближения) и с внешним R-пакетом
\texttt{npSPC}; статистики и сигналы совпадают, а расхождения чужого кода
документированы и не переносятся в нашу реализацию. Единственное
открытое расхождение --- рис.~4.6 книги --- зафиксировано без подгонки.
Практическая рекомендация: использовать консервативный подбор $h$ и
вычисленные дизайны $(k,h)$, а не печатные значения рис.~4.6.

% ================================================================== %
\appendix
\section{Приложения}
\label{sec:appendix}
% ================================================================== %

\paragraph{Воспроизведение.} Среда описана в
\path{experiments/00_ENVIRONMENT.txt}; восстановление R-окружения
(микромamba, R, \texttt{npSPC}) --- скриптом
\path{experiments/cusum/setup_renv.sh}. Порядок запуска:
\begin{lstlisting}[language=bash]
poetry run python experiments/cusum/01_oc_curves.py        # OC curves
poetry run python experiments/cusum/02_sr_exact_vs_mc.py   # quadrature vs MC
poetry run python experiments/cusum/03_kh_design.py        # (k, h) design
poetry run python experiments/cusum/04_performance.py      # performance
RSCRIPT=$RENV/bin/Rscript R_LIBS=$RENV/lib/R/library \
    poetry run python experiments/cusum/05_cross_validation.py
poetry run python experiments/cusum/06_gen_figdata.py \
    > docs/reports/cusum_macros.tex
cd docs/reports && make cusum
\end{lstlisting}
Тесты: \texttt{poetry run pytest -q} (включая медленные
\texttt{test\_run\_length\_performance.py}; быстрый профиль ---
\texttt{npc test --fast}). Скрипты 03 и 04 печатают таблицы в консоль и
одновременно сохраняют их в \path{results/experiments/cusum_kh_design.txt}
и \path{results/experiments/cusum_performance.txt}; сверка Python$\leftrightarrow$R
сохраняется в \path{results/experiments/cusum_cross_validation.txt}.

\paragraph{Источники формул.} Уравнения рекурсии CUSUM-SN --- 4.7--4.8
(стр.~203); CUSUM-SR --- 4.10--4.11 (стр.~206, с опечаткой 4.11,
таблица~\ref{tab:typo}); матрица переходов и PMF --- 4.9 и 4.12
(стр.~203, 206), уравнения длины серии --- 3.22--3.24 (стр.~112);
дизайн --- \S3.2.2 (стр.~67--72). Первоисточник непараметрического
CUSUM-SR --- Bakir и Reynolds (1979).

% ================================================================== %
\begin{thebibliography}{9}
% ================================================================== %

\bibitem{Amin1995}
Amin, R.~W., Reynolds~Jr., M.~R., Bakir, S.~T. (1995).
Nonparametric quality control charts based on the sign statistic.
\emph{Communications in Statistics: Theory and Methods}, 24(6), 1597--1623.

\bibitem{Bakir1979}
Bakir, S.~T., Reynolds~Jr., M.~R. (1979).
A nonparametric procedure for process control based on within-group ranking.
\emph{Technometrics}, 21(2), 175--183.

\bibitem{Chakraborti2019}
Chakraborti, S., Graham, M.~A. (2019).
\emph{Nonparametric Statistical Process Control}. Wiley.

\bibitem{Siegmund1985}
Siegmund, D. (1985).
\emph{Sequential Analysis: Tests and Confidence Intervals}. Springer.

\bibitem{Nascimento2026}
Nascimento, D.~C., Fons\^eca, A.~O., Fons\^eca, M.~B., Silva, M.~A.~F.,
Ferreira, P.~H. (2026).
Nonparametric SPC: Solving complex industrial monitoring tasks with low
computational cost. \emph{Array}, 30, 100828.
\url{https://doi.org/10.1016/j.array.2026.100828}.

\bibitem{Flores2021}
Flores, M., Fern\'andez-Casal, R., Naya, S., Tarr\'io-Saavedra, J. (2021).
qcr: Quality Control Review. \emph{The R Journal}.
\url{https://doi.org/10.32614/RJ-2021-034}.

\end{thebibliography}

\end{document}
